\section{Introduction}
\label{sec:introduction}

When language models train iteratively on their own output, the training distribution progressively loses its tail modes, a phenomenon termed model collapse~\cite{shumailov2024curse,alemohammad2024mad}.
Lexical diversity narrows, calibration deteriorates, and each generation inherits the systematic biases of its predecessors.
As machine-generated text becomes prevalent on the web, such feedback loops pose a growing and often invisible risk to production training pipelines.
By the time standard quality metrics register degradation, the underlying distribution may have already shifted irreversibly.
This threat motivates the search for cheap-to-compute \emph{canary metrics}: leading indicators that signal distributional degradation before expensive downstream measures reveal the full extent of collapse.
The concept parallels early-warning signals for critical transitions in ecological systems~\cite{scheffer2009early,scheffer2012anticipating}, where proxy indicators allow intervention before a tipping point.

Yet a fundamental question remains unanswered: \emph{at what contamination level does Granger-causal detection of collapse become possible?}
Whether a candidate canary genuinely \emph{precedes} downstream degradation or merely correlates with it under shared non-stationary trends depends on contamination regime and statistical method in ways that existing analyses leave unresolved.
Iterative training produces short, noisy time series (often fewer than 15 points per generation sequence), where detection power drops sharply and false positive rates can be high~\cite{boettiger2012limits}.

We identify four deficiencies in existing lead-lag analyses of iterative training collapse.
(1)~No differencing tests rule out spurious regression from shared non-stationary trends~\cite{granger1974spurious}.
(2)~Conclusions rest on a single statistical method whose assumptions about stationarity and directionality are rarely examined.
(3)~No FDR correction guards against spurious findings when testing many metric pairs simultaneously.
(4)~No regime-dependent analysis is performed, despite the possibility that the same test yields opposite conclusions at different contamination levels.

We discover a \emph{sharp detection onset} in canary-based monitoring of iterative collapse.
Using Pythia-410M~\cite{biderman2023pythia} across thirteen contamination levels ($\alpha \in \{0.00, 0.05, 0.10, 0.25, 0.30, 0.35, 0.40, 0.45, 0.50, 0.55, 0.60, 0.75, 1.00\}$) and up to four random seeds, we show that causal detectability undergoes a steep onset at per-generation mixing ratio $\alpha_c{=}0.476$ at $T{=}10$ generations (sigmoid $R^2{=}0.969$, $\Delta\alpha{\approx}0.02$): permutation Granger causality (PermG) jumps from 0/8 FDR-significant metric pairs at $\alpha{\leq}0.45$ to 6/8 at $\alpha{=}0.50$.
A signal-times-resolution ($S \times R$) framework explains this onset and the non-monotonic dose-response through competing signal strength and temporal resolution.
A multi-method consensus framework (MMCT) comparing five statistical methods reveals that method choice alone can shift conclusions on the same data from universal significance to complete null.
Cross-architecture validation across 15 model-seed combinations spanning six families (345M--6.9B) confirms PermG detection (${\geq}6/8$) in 13 cases (87\%).
Granger precedence in our framework establishes temporal predictive improvement, not true causality~\cite{shojaie2022granger}; our canary claim is scoped to diversity metrics (distinct-$n$~\cite{li2016diversity}), not model collapse broadly.

Our contributions are:
\begin{enumerate}
    \item \textbf{Sharp detection onset} (\S\ref{sec:regime}).
    At $T{=}10$, PermG jumps from 0/8 FDR-significant pairs through $\alpha{=}0.45$ to 6/8 at $\alpha{=}0.50$ ($\alpha_c{=}0.476$, sigmoid $R^2{=}0.969$, $\Delta\alpha{\approx}0.02$), with a four-regime structure: \emph{blind}, \emph{onset}, \emph{synchronous dip}, and \emph{detection}.
    At $\alpha{=}0.10$, a canary blind spot emerges: diversity has already collapsed but canary metrics remain indistinguishable from the null.
    Cross-architecture validation across 15 model-seed combinations confirms detection in 13 cases (87\%).

    \item \textbf{Signal-times-resolution theory} (\S\ref{sec:theory}).
    A power interpretation formalizes the onset as $\mathrm{Power}(\alpha, T) = \Phi(S \cdot R \cdot \sqrt{T}/\sigma)$, generating three testable predictions for $T{=}20$; sub-threshold PermG $p$-values at $\alpha{=}0.45$ ($p{=}0.059$--$0.091$) are consistent with proximity to the power boundary.

    \item \textbf{Multi-method diagnostic toolkit} (\S\ref{sec:method}).
    The MMCT exposes three pitfalls: method choice alone inflates spurious-signal rates to 75\% ($\kappa{=}0.187$), first-differencing eliminates canary$\to$perplexity associations entirely while preserving all canary$\to$diversity leads (24/24$\to$24/24), and regime dependence causes the same test to shift from 5--6/8 FDR-significant pairs per seed at $\alpha{\geq}0.75$ to 0/8 at $\alpha{\leq}0.45$.
    PermG identifies token entropy as a one-generation-ahead canary for diversity collapse (24/24 forward, 0/24 reverse).

    \item \textbf{Canary-based early detection} (\S\ref{sec:intervention}).
    Predictive cessation reduces peak perplexity by 58\% (bootstrap 95\% CI: 29--64\%) and fires one generation earlier than reactive cessation.
    A four-layer detection taxonomy provides regime-appropriate diagnostics from $\alpha{\geq}0.10$ to $\alpha{\geq}0.50$.
\end{enumerate}
